
% ============================================================================
% BESOINS NON-FONCTIONNELS
% MOUTENG REBECCA ESTHER - Matricule: 25G2079
% ============================================================================

\chapter{Exigences Non-Fonctionnelles et Qualité}

\section{Attributs de Qualité}

%\textit{À COMPLÉTER PAR MOUTENG}
La qualité d'un système logiciel ne se limite pas à sa capacité à réaliser les fonctionnalités attendues. Elle s'évalue
également à travers un ensemble de propriétés non fonctionnelles
qui déterminent comment le système se comporte dans son environnement
réel d'utilisation.

% Consultez le fichier MOUTENG/task.tex pour les ressources complètes
% Couvrir : Sécurité, Disponibilité, Performance, Maintenabilité, Scalabilité, Usabilité

\subsection{Sécurité}



La sécurité protège les données pédagogiques et personnelles des utilisateurs.

\begin{description}

    \item[ — Authentification:]
    Accès par identifiant et mot de passe obligatoire. Mots de passe
    hachés en \textbf{bcrypt/Argon2} (12 rounds minimum). Compte verrouillé
    \textbf{30 minutes} après \textbf{5 échecs} consécutifs.\\
    \textit{Métrique :} 0 mot de passe en clair ;
    verrouillage après $\leq 5$ tentatives.\\
    \textit{Justification :} Empêcher l'accès libre et les attaques
    par force brute.

    \item[ — Contrôle d'accès RBAC:]
    4 rôles définis : Enseignant, Étudiant, Parent, Administrateur.
    Un étudiant ne voit que sa classe, un parent que ses enfants,
    un enseignant ne modifie que ses séances.\\
    \textit{Métrique :} 0 accès inter-rôle non autorisé ;
    tentative illégale retourne HTTP \texttt{403}.\\
    \textit{Justification :} Principe du moindre privilège.

    \item[ — HTTPS obligatoire:]
    Toutes les communications via \textbf{HTTPS} (TLS 1.2 minimum).
    MD5 et SHA-1 strictement interdits.\\
    \textit{Métrique :} 0 endpoint HTTP non chiffré en production.\\
    \textit{Justification :} Prévenir l'interception des données
    sur les réseaux du campus.

    \item[ — Journalisation et immuabilité des logs:]
    Tous les accès et modifications enregistrés (qui, quand, quoi),
    conservés \textbf{1 an minimum}. Suppression des logs interdite.\\
    \textit{Métrique :} Couverture 100\% ; logs $\geq 365$ jours ;
    0 droit \texttt{DELETE}/\texttt{UPDATE} sur la table des logs.\\
    \textit{Justification :} Garantir la valeur probatoire des journaux
    en cas d'incident.

\end{description}

\subsection{Disponibilité}


Elle garantit que le système est accessible aux utilisateurs
au moment où ils en ont besoin.

\begin{description}

    \item[ — Uptime minimum garanti:]
    Le système \textbf{DOIT} être disponible au minimum \textbf{99\%}
    du temps (max 7,2h de panne/mois), avec monitoring 24/7
    et alertes automatiques.\\
    \textit{Métrique :} Uptime mensuel $\geq 99\%$.\\
    \textit{Justification :} Garantir l'accès pendant les heures
    d'activité universitaire.

    \item[ — Fenêtres de maintenance contrôlées:]
    Maintenance programmée le dimanche entre \textbf{22h et 00h}.
    Patch critique déployé en moins de \textbf{30 minutes}.\\
    \textit{Métrique :} 100\% des maintenances dans la fenêtre
    autorisée ; patch critique $\leq 30$ min.\\
    \textit{Justification :} Minimiser l'impact sur les utilisateurs.

    \item[ — Architecture redondante:]
    Base de données avec replica en standby (failover automatique).
    Serveur applicatif sur \textbf{2 instances minimum} derrière
    un load balancer.\\
    \textit{Métrique :} Failover $\leq 60$s ;
    0 point de défaillance unique (SPOF).\\
    \textit{Justification :} Éviter qu'une panne d'un composant
    rende le système inaccessible.

    \item[ — Reprise après sinistre:]
    RTO $\leq$ \textbf{4h}, RPO $\leq$ \textbf{24h}.
    Backup quotidien stocké hors-site.\\
    \textit{Métrique :} RTO $\leq 4$h ; RPO $\leq 24$h ;
    backup quotidien vérifié.\\
    \textit{Justification :} Protéger les données pédagogiques
    et assurer une reprise rapide.

\end{description}
% -----------------------------------------------------------------------

\subsection{Performance}


Elle s'assure que le système répond rapidement aux actions des
utilisateurs.

\begin{description}

    \item[ — Temps de réponse par opération.]
    Les temps de réponse au \textbf{95\textsuperscript{e} percentile} sont :
    login $< 1$s, consultation $< 2$s, saisie de séance $< 3$s,
    génération de rapport $< 10$s.\\
    \textit{Métrique :} 95\% des requêtes respectent ces seuils
    (mesuré par monitoring en production).\\
    \textit{Justification :} Garantir une expérience fluide pour
    les enseignants et étudiants.

    \item[ — Charge simultanée supportée.]
    Le système \textbf{DOIT} supporter \textbf{1000 requêtes/minute}
    et \textbf{500 utilisateurs simultanés} sans dégradation.\\
    \textit{Métrique :} Temps de réponse maintenus sous 500 utilisateurs
    simultanés (vérifié par tests de charge).\\
    \textit{Justification :} Absorber les pics en début et fin de
    journée universitaire.

    \item[ — Cache multi-niveaux.]
    Cache navigateur (données statiques : \textbf{1 jour}), cache
    \textbf{Redis} (données dynamiques : \textbf{1 heure}), index base
    de données sur colonnes fréquemment requêtées.\\
    \textit{Métrique :} Taux de cache hit $\geq 70\%$ sur données dynamiques.\\
    \textit{Justification :} Réduire les appels répétitifs à la
    base de données.

    \item[ — Ressources serveur maîtrisées.]
    CPU $\leq 70\%$ et RAM $\leq 80\%$ en charge normale. Tests de
    performance \textbf{obligatoires} avant tout déploiement.\\
    \textit{Métrique :} Seuils respectés sous 500 utilisateurs simultanés.\\
    \textit{Justification :} Maintenir une marge pour absorber
    les pics inattendus.

\end{description}



\subsection{Maintenabilité}



 Garantit que le système peut être facilement compris,
corrigé et faire l'objet d'évolutions sans effort disproportionné.

\begin{description}

    \item[ — Modularité du code.]
    Le code \textbf{DOIT} être organisé en modules indépendants, faiblement
    couplés et fortement cohésifs. Chaque module ne \textbf{DOIT} pas
    dépasser \textbf{500 lignes} ou \textbf{20 classes}.\\
    \textit{Métrique :} 0 module dépassant les seuils définis
    (vérifié par analyse statique du code).\\
    \textit{Justification :} Faciliter la localisation et la correction
    des anomalies sans effets de bord.

    \item[ — Documentation.]
    Le code \textbf{DOIT} être commenté avec un ratio minimum de
    \textbf{30\% commentaires/code}. L'architecture, les APIs et les
    procédures de déploiement \textbf{DOIVENT} être documentées
    (README, Swagger/OpenAPI).\\
    \textit{Métrique :} Ratio commentaires $\geq 30\%$ ;
    documentation API à jour à chaque release.\\
    \textit{Justification :} Permettre à tout nouveau membre de l'équipe
    de comprendre et modifier le système rapidement.

    \item[ — Couverture de tests.]
    La couverture minimale par les tests unitaires \textbf{DOIT} être de
    \textbf{80\%}. Les modules critiques \textbf{DOIVENT} disposer de tests
    d'intégration. Des tests de performance \textbf{DOIVENT} être exécutés
    avant chaque déploiement.\\
    \textit{Métrique :} Couverture tests unitaires $\geq 80\%$
    (mesuré par outil de couverture de code).\\
    \textit{Justification :} Détecter les régressions tôt et réduire
    le risque lors des mises à jour.

    \item[ — Versioning et traçabilité.]
    Le code source \textbf{DOIT} être géré sous \textbf{Git} avec des
    commits atomiques et des messages clairs. Le schéma de base de données
    et les releases \textbf{DOIVENT} suivre le versioning sémantique
    (\texttt{major.minor.patch}).\\
    \textit{Métrique :} 100\% du code source sous Git ;
    0 commit sans message descriptif.\\
    \textit{Justification :} Permettre le retour arrière rapide en cas
    de régression et faciliter la collaboration entre développeurs.

\end{description}

\subsection{Scalabilité}



La scalabilité garantit que le système peut supporter une croissance
du nombre d'utilisateurs et de données sans modification majeure
de l'architecture.

\begin{description}

    \item[ — Scalabilité horizontale.]
    L'architecture \textbf{DOIT} être multi-tiers, stateless et permettre
    l'ajout de serveurs via un \textbf{load balancer} sans interruption
    de service.\\
    \textit{Métrique :} Ajout d'une instance serveur sans downtime ;
    0 dépendance de session côté serveur.\\
    \textit{Justification :} Absorber la croissance du nombre d'utilisateurs
    par simple ajout de machines.

    \item[ — Scalabilité verticale.]
    La base de données \textbf{DOIT} supporter \textbf{10 000+ séances}
    avec des performances maintenues grâce aux index et au partitioning.\\
    \textit{Métrique :} Temps de requête $< 2$s jusqu'à 10 000 séances
    (vérifié par tests de charge).\\
    \textit{Justification :} Éviter toute dégradation au fur et à mesure
    que le volume de données augmente.

    \item[ — Phases de croissance supportées.]
    Le système \textbf{DOIT} supporter trois phases :
    \textbf{Phase 1} (100 utilisateurs — pilote),
    \textbf{Phase 2} (1000 utilisateurs — déploiement complet),
    \textbf{Phase 3} (10 000 utilisateurs — extension multi-sites).\\
    \textit{Métrique :} Performances de NFR-PERF-01 respectées
    à chaque phase.\\
    \textit{Justification :} Anticiper la croissance progressive
    à l'échelle de l'université.

    \item[ — Indépendance des composants.]
    Chaque composant (authentification, cahier de texte, rapports)
    \textbf{DOIT} pouvoir être mis à l'échelle indépendamment des autres.\\
    \textit{Métrique :} 0 couplage fort entre composants empêchant
    un scaling indépendant.\\
    \textit{Justification :} Optimiser les ressources en ne scalant
    que les composants sous forte charge.

\end{description}

\subsection{Usabilité}

Elle permet  que le cahier de texte digital est simple et agréable
à utiliser pour tous les acteurs.

\begin{description}

    \item[ — Interface intuitive.]
    L'interface \textbf{DOIT} être épurée, avec des menus clairs et des
    boutons explicitement étiquetés (ex. ``Créer séance'' et non
    ``Ajouter''). Aucune formation préalable ne doit être nécessaire
    pour les tâches courantes.\\
    \textit{Métrique :} Taux de réussite des tâches courantes $\geq 90\%$
    lors des tests utilisateurs.\\
    \textit{Justification :} Favoriser l'adoption rapide par les enseignants
    et personnels administratifs.

    \item[ — Accessibilité.]
    L'interface \textbf{DOIT} respecter les normes \textbf{WCAG 2.1 niveau AA} :
    contraste minimum \textbf{4.5:1}, navigation complète au clavier,
    support des lecteurs d'écran.\\
    \textit{Métrique :} 0 violation WCAG 2.1 AA
    (vérifié par outil d'audit d'accessibilité).\\
    \textit{Justification :} Garantir l'accès au système aux utilisateurs
    en situation de handicap.

    \item[ — Responsive design.]
    L'interface \textbf{DOIT} fonctionner correctement sur desktop, tablette
    et mobile avec les breakpoints : \textbf{320px} (mobile),
    \textbf{768px} (tablette), \textbf{1024px} (desktop).\\
    \textit{Métrique :} 0 dysfonctionnement d'affichage sur les trois
    breakpoints (vérifié par tests multi-supports).\\
    \textit{Justification :} Permettre aux enseignants de saisir leurs
    séances depuis n'importe quel appareil.

    \item[ — Localisation.]
    L'interface \textbf{DOIT} être disponible en \textbf{français}.
    Une extension vers l'anglais \textbf{DOIT} être possible sans
    modification de l'architecture.\\
    \textit{Métrique :} 100\% des libellés en français ; ajout d'une
    langue sans refactoring du code.\\
    \textit{Justification :} Répondre au contexte bilingue de
    l'Université de Yaoundé I.

\end{description}



\section{Tableau Récapitulatif des Exigences}


\begin{table}[H]
\centering
\caption{Récapitulatif des Exigences Non-Fonctionnelles}
\label{tab:nfr}
\small
\begin{tabularx}{\textwidth}{|l|X|X|l|}
\hline
\textbf{Attribut} & \textbf{Description}
& \textbf{Métrique} & \textbf{Priorité} \\
\hline

% --- SÉCURITÉ ---
Sécurité
& Authentification, hachage bcrypt/Argon2, verrouillage après 5 échecs
& 0 mdp en clair ; blocage $\leq 5$ tentatives & Critique \\
\hline
Sécurité
& Contrôle d'accès RBAC (4 rôles)
& 0 accès inter-rôle non autorisé & Critique \\
\hline
Sécurité
& HTTPS obligatoire (TLS 1.2+)
& 0 endpoint HTTP non chiffré & Critique \\
\hline
Sécurité
& Journalisation et immuabilité des logs
& Couverture 100\% ; logs $\geq 365$ jours & Critique \\
\hline

% --- DISPONIBILITÉ ---
Disponibilité
& Uptime minimum garanti
& Uptime $\geq 99\%$ & Critique \\
\hline
Disponibilité
& Fenêtres de maintenance contrôlées
& Patch critique $\leq 30$ min & Élevée \\
\hline
Disponibilité
& Architecture redondante
& Failover $\leq 60$s ; 0 SPOF & Critique \\
\hline
Disponibilité
& Reprise après sinistre
& RTO $\leq 4$h ; RPO $\leq 24$h & Élevée \\
\hline

% --- PERFORMANCE ---
Performance
& Temps de réponse par opération
& 95\% requêtes dans les seuils & Élevée \\
\hline
Performance
& Charge simultanée supportée
& 500 utilisateurs sans dégradation & Élevée \\
\hline
Performance
& Cache multi-niveaux
& Cache hit $\geq 70\%$ & Moyenne \\
\hline
Performance
& Ressources serveur maîtrisées
& CPU $\leq 70\%$ ; RAM $\leq 80\%$ & Moyenne \\
\hline

% --- MAINTENABILITÉ ---
Maintenabilité
& Modularité du code
& 0 module $>$ 500 lignes ou 20 classes & Élevée \\
\hline
Maintenabilité
& Documentation
& Ratio commentaires $\geq 30\%$ & Élevée \\
\hline
Maintenabilité
& Couverture de tests
& Tests unitaires $\geq 80\%$ & Élevée \\
\hline
Maintenabilité
& Versioning Git
& 100\% code sous Git & Moyenne \\
\hline

% --- SCALABILITÉ ---
Scalabilité
& Architecture stateless multi-tiers
& 0 dépendance session côté serveur & Élevée \\
\hline
Scalabilité
& Base de données optimisée
& Requêtes $< 2$s jusqu'à 10 000 séances & Élevée \\
\hline
Scalabilité
& Phases de croissance supportées
& Performances maintenues à chaque phase & Moyenne \\
\hline
Scalabilité
& Indépendance des composants
& 0 couplage fort entre composants & Moyenne \\
\hline

% --- USABILITÉ ---
Usabilité
& Interface intuitive
& Taux de réussite tâches $\geq 90\%$ & Élevée \\
\hline
Usabilité
& Accessibilité WCAG 2.1 AA
& 0 violation WCAG 2.1 AA & Élevée \\
\hline
Usabilité
& Responsive design
& 0 bug sur les 3 breakpoints & Élevée \\
\hline
Usabilité
& Localisation français/anglais
& Ajout de langue sans refactoring & Moyenne \\
\hline

\end{tabularx}
\end{table}

\section{Contraintes}

La conception du cahier de texte digital ne se limite pas aux seuls attributs de qualité. Elle s'inscrit dans
un cadre plus large de contraintes qui s'imposent au système
indépendamment des choix de l'équipe de développement.Elles
sont de trois natures : technologiques, réglementaires et budgétaires.

\subsection{Contraintes Technologiques}


Elles définissent le cadre technique dans lequel
le système doit être développé et déployé.

\begin{description}

    \item[CT-01 — Langages et frameworks.]
    Le système \textbf{DOIT} être développé avec des technologies web
    standard (HTML, CSS, JavaScript) et un framework backend reconnu
    (ex. Spring Boot, Laravel, Django). L'utilisation de technologies
    obsolètes est interdite.\\
    \textit{Justification :} Garantir la maintenabilité et la disponibilité
    de compétences sur le marché local.

    \item[CT-02 — Base de données.]
    Le système \textbf{DOIT} utiliser un SGBD relationnel
    (PostgreSQL ou MySQL). La base \textbf{DOIT} supporter les transactions
    ACID et les mécanismes de réplication.\\
    \textit{Justification :} Assurer l'intégrité des données pédagogiques
    et la disponibilité du système.

    \item[CT-03 — Déploiement.]
    Le système \textbf{DOIT} être déployable sur un serveur Linux
    (Ubuntu 20.04+). La conteneurisation via \textbf{Docker} est recommandée
    pour faciliter les déploiements et la portabilité.\\
    \textit{Justification :} Réduire les dépendances à un environnement
    spécifique et simplifier la maintenance.

    \item[CT-04 — Navigateurs supportés.]
    L'interface \textbf{DOIT} fonctionner sur les versions récentes de
    Chrome, Firefox, Edge et Safari. Le support d'Internet Explorer
    est exclu.\\
    \textit{Justification :} Couvrir l'ensemble des navigateurs utilisés
    sur le campus de Yaoundé I.

\end{description}

\subsection{Contraintes Réglementaires}
Elles mettent en avant les obligations légales
auxquelles le système doit se conformer.

\begin{description}

    \item[CR-01 — Protection des données personnelles.]
    Le système \textbf{DOIT} respecter les principes du RGPD : collecte
    minimale des données, consentement explicite, droit d'accès et
    droit à l'oubli.\\
    \textit{Justification :} Protéger les données personnelles des
    étudiants, enseignants et parents conformément aux réglementations
    en vigueur.

    \item[CR-02 — Confidentialité des données pédagogiques.]
    Les données pédagogiques (notes, présences, séances) sont
    confidentielles et ne \textbf{DOIVENT} être accessibles qu'aux
    personnes autorisées selon leur rôle.\\
    \textit{Justification :} Respecter le secret professionnel et les
    règlements intérieurs de l'Université de Yaoundé I.

    \item[CR-03 — Conservation des données.]
    Les données pédagogiques \textbf{DOIVENT} être conservées au minimum
    \textbf{5 ans} conformément aux obligations légales des établissements
    d'enseignement supérieur.\\
    \textit{Justification :} Permettre la consultation des archives en
    cas de litige ou de vérification administrative.

\end{description}



\subsection{Contraintes de Coûts}

 Cette catégorie de besoins non fonctionnels permet de définir  les limites budgétaires dans
lesquelles le système doit être conçu et maintenu.

\begin{description}

    \item[CC-01 — Technologies open source.]
    Le système \textbf{DOIT} privilégier les technologies open source
    (Linux, PostgreSQL, frameworks libres) afin de réduire les
    coûts de licence.\\
    \textit{Justification :} Adapter le projet au budget limité d'un
    établissement universitaire public.

    \item[CC-02 — Hébergement.]
    Le système \textbf{DOIT} pouvoir être hébergé sur l'infrastructure
    existante de l'université ou sur un serveur cloud économique
    (ex. VPS).\\
    \textit{Justification :} Éviter des investissements matériels
    importants en phase initiale.

    \item[CC-03 — Maintenance.]
    Le coût de maintenance annuel \textbf{NE DOIT PAS} dépasser le budget
    alloué par l'université. La documentation et les tests \textbf{DOIVENT}
    réduire le temps d'intervention en cas d'anomalie.\\
    \textit{Justification :} Assurer la pérennité du système sur le long
    terme sans surcharge budgétaire.

\end{description}

\section{Trade-offs}


Les trade-offs identifient les tensions entre exigences qui nécessitent
des compromis lors de la conception et du développement du système.

\begin{description}

    \item[TF01 — Sécurité vs Performance.]
    Le chiffrement HTTPS, le hachage bcrypt et la journalisation exhaustive
    introduisent une latence supplémentaire pouvant impacter les temps de
    réponse définis.\\
    \textit{Compromis retenu :} La sécurité est prioritaire. Les performances
    seront optimisées via le cache pour compenser la latence
    introduite par le chiffrement.

    \item[TF-02 — Disponibilité vs Coûts.]
    Une architecture redondante  avec load balancer et replicas
    de base de données implique des coûts d'infrastructure élevés, en tension
    avec les contraintes budgétaires.\\
    \textit{Compromis retenu :} Déploiement progressif — redondance minimale
    en Phase 1 (pilote), architecture complète en Phase 2 (déploiement full).

    \item[TF-03 — Performance vs Maintenabilité.]
    L'optimisation poussée des performances 
    produit un code plus complexe.\\
    \textit{Compromis retenu :} Optimiser uniquement les modules critiques
    identifiés par profiling. Le reste du code privilégie la lisibilité.


    \item[TF-04 — Sécurité vs Usabilité.]
    Le verrouillage de compte après 5 échecs  et les restrictions peuvent frustrer les utilisateurs.\\
    \textit{Compromis retenu :} Mettre en place un mécanisme de déverrouillage
    par email et des messages d'erreur clairs guidant l'utilisateur sans
    révéler d'informations sensibles.

\end{description}

\vfill
